\section{Discussion}
\label{sec:discussion}

\subsection{Interpretation of Estimates}

Across the boundary and redistricting designs, more stringent aldermen reduce housing supply on both the density and permitting margins.
The rent and sale-price comparisons suggest that these supply differences may also be reflected in local housing costs, an important concern as Chicago rents continue to rise rapidly, outpacing much of the country \citep{davis_chicagos_2026}.
Aldermanic privilege creates sharp changes in the regulatory environment at ward boundaries.
Those inconsistencies make the development process slower and less predictable, and likely deter some projects from ever breaking ground.

\citet{khan_decentralized_2021} documents the intensive margin of this institution: conditional on a developer seeking a rezoning, 
aldermen grant smaller rezonings on average near ward boundaries where they cannot internalize the positive benefits of development. 
My density estimates capture that margin, but they also speak to an broader margin of realized development, since allowed density and realized built density can differ substantially, especially in the presence of discretionary review \cite{rollet_zoning_2025}.
What neither my estimates nor Khan's can capture is the deterrence margin: projects that are never proposed because developers anticipate rejection or sort into more friendly locations elsewhere.
The Fern Hill episode described in Section~\ref{sec:data} illustrates both observed margins.
The approved project was smaller than originally proposed, and four years of community meetings and aldermanic negotiation signal to other developers what attempting a large project can entail.
The time-to-build framework in Section~\ref{sec:model} helps explain why deterrence may be quantitatively important: even a modest increase in expected regulatory delay can flip marginal projects from viable to unviable, and those projects leave no trace in the data.
My estimates are therefore best read as lower bounds on the total supply effect of aldermanic discretion.

\subsection{Policy Implications}

The results speak directly to current debates about reforming aldermanic privilege, which have become a hot topic in Chicago politics in recent years.
The supply effects I document build through longer processing times, greater uncertainty, and the resulting deterrence of marginal projects, suggesting multiple cascading consequences on the 
Chicago housing market and the spatial distribution of development across the city.

Centralizing land use authority by shifting discretion from individual aldermen to citywide bodies would let decision-makers internalize the full benefits of housing production, though at the cost of sacrificing some local knowledge.
Expanding by-right development so that code-compliant projects proceed automatically without aldermanic review would remove the discretionary layer while preserving community preferences as expressed in the zoning code.
But as Section~\ref{sec:data} shows, aldermen retain unilateral power to downzone parcels preemptively, which means they can often undo by-right protections before a developer applies.
The Reilly and Solis episodes make the point clearly: reforms that target only the permitting process, without constraining aldermanic control over the zoning code itself, may have limited effect.
Effective reform likely needs to address both channels.
